\section*{Materials and methods}

% Owner: papers/1_method/archive/section-plans-20260729/methods.md (M1--M11) + papers/_program/provenance.md (run manifest).
% Numbers are annotated with `% src:` per LINT-SRC. Definitions owned by the Theory
% and Diagnostics sections are cited here, not restated (papers/_program/00_index.md).

\subsection*{Minimal model and simulation}

We work with a two-state channel scheme, closed ($C$) and open ($O$), the smallest
scheme in which the acquisition-averaging effect we study is present.
% src: papers/1_method/archive/section-plans-20260729/methods.md#M1 (scheme_CO; legacy/models_simple.h:10-66)
The closed-to-open transition is agonist dependent with association constant
$k_{\mathrm{on}}$ multiplying the agonist concentration; the open-to-closed
transition is agonist independent with dissociation rate $k_{\mathrm{off}}$. Only
the open state conducts, and the code applies a sign flip so the single-channel
current is inward. At the start of every recording all channels are closed, so the
occupancy is $(1,0)$.
% src: papers/1_method/archive/section-plans-20260729/methods.md#M1 (v_g_formula; P_initial=[1,0]; sign flip)

The scheme carries six parameters, all estimated in base-ten logarithmic
coordinates. The values used to generate the data reported here, read from the
production run scripts rather than from the parameter tables in the repository
(the tables hold stale values the figures never used), are
$k_{\mathrm{on}} = 10\;\mu\mathrm{M}^{-1}\,\mathrm{s}^{-1}$,
$k_{\mathrm{off}} = 100\;\mathrm{s}^{-1}$,
a unitary conductance $g$ set to $1$ (the engine parameter
\texttt{unitary\_current}, written $i$ where the noise axis is defined below), and a
current baseline $b$ set to $1$ (both in units of $g$),
% src: projects/eLife_2025/ops/local/figure_3_mle.macroir:42-50 (on=10, off=100, unitary=1, baseline=1)
% src: papers/1_method/decisions/D-2_parameter_units.md (settled units; CSV != run values)
while the white-noise level and the mean channel number were swept (below). Current
is expressed throughout in units of the unitary conductance $g$, which corresponds
physically to a single-channel current of roughly $1\;\mathrm{pA}$.
% src: papers/1_method/decisions/D-2_parameter_units.md (current in units of g; g ~ 1 pA)
The dissociation time $\tau = 1/k_{\mathrm{off}} = 10\;\mathrm{ms}$ sets the natural
time scale that normalises the acquisition-interval axis. The relaxation at the
agonist condition is faster, with rate $k_{\mathrm{on}}[A] + k_{\mathrm{off}} =
200\;\mathrm{s}^{-1}$; the axis uses $\tau = 1/k_{\mathrm{off}}$ by convention.
% src: papers/1_method/decisions/D-2_parameter_units.md (tau = 1/k_off = 10 ms; lambda = k_on[A]+k_off)

The emission model is white instrumental noise only. Both the pink (one-over-$f$)
and the proportional noise channels are set to zero, so the term the reader may
expect from flicker noise is genuinely absent.
% src: papers/1_method/archive/section-plans-20260729/methods.md#M1 (pink and proportional = 0)
For a measurement interval $t$ of duration $\Delta_t$, the predicted current has
mean $\bar y_t = N_{\mathrm{ch}}\,\mu_t\!\cdot\!\gamma_t + b$ and variance
\[
  \sigma_t^2 \;=\; e_t \;+\; N_{\mathrm{ch}}\,v_t ,
  \qquad
  e_t \;=\; \mathrm{Current\_Noise}\,\frac{f_s}{n_{\mathrm{samp},t}}
        \;=\; \frac{\mathrm{Current\_Noise}}{\Delta_t},
\]
where $N_{\mathrm{ch}}$ is the channel number, $\mu_t$ the occupancy mean, $\gamma_t$
the interval-averaged conductance, $f_s$ the raw sampling rate, $n_{\mathrm{samp},t}$
the number of raw samples averaged into the interval, $e_t$ the white instrumental
variance over the interval, and $N_{\mathrm{ch}}\,v_t$ the channel-gating variance.
% src: papers/1_method/decisions/D-2_parameter_units.md (e = Current_Noise*fs/n_samp)
% src: legacy/qmodel.h:3740-3741 (e = Current_Noise * fs / number_of_samples + Pink_Noise)
The white variance $e_t$ falls as the interval lengthens, because more raw samples
are averaged into each recorded value. The gating variance $v_t$ is the
single-interval conductance variance, whose form differs across the algorithm
family and is derived in the Theory section.

The stimulus is a single concentration jump in three segments, a pre-pulse at
$0\;\mu\mathrm{M}$, an agonist step at $10\;\mu\mathrm{M}$, and a post-pulse at
$0\;\mu\mathrm{M}$, with raw sampling at $f_s = 50\;\mathrm{kHz}$ throughout.
% src: projects/eLife_2025/ops/local/figure_3_mle.macroir:54-61 (three segments; 0/10/0 uM; fs=50e3)
Each recorded measurement is the mean of $n_{\mathrm{samp},t}$ consecutive raw
samples, so the averaging window is fixed at acquisition. That window is the
acquisition interval the whole study is about, and its length is set by the
number of raw samples entering each recorded value.
% src: papers/1_method/archive/section-plans-20260729/methods.md#M2 (interval = averaging window, not a filter)
The averaging is uniform over the window and no analogue filter stage is simulated,
so a recorded value is a boxcar average of the trajectory, which is the observable
every likelihood in the family is written for (Theory). Recordings from a real
amplifier are Bessel filtered instead, and that kernel is outside the scope of this
paper (Discussion).
% [BESSEL-BOUND] OPEN 2026-08-06. Canonical note at 02_theory.tex, under the Bessel paragraph;
% the decision on how much to print is at 05_discussion.tex, beside the Bessel paragraph there.
% THIS IS WHERE A READER WITH A RECORDING LOOKS, and the sentence above sends them away without
% an instruction. "Outside the scope of this paper" is true and unhelpful: it does not say
% whether their record can be brought inside the scope, and it can, by block-averaging.
% THE OPERATIONAL FORM, which costs about 40 words and no run:
%   "A record from an amplifier can be brought to this observable by averaging consecutive
%   samples into windows long against the filter's response time, since the average of an
%   average is again an average over the longer window; what remains unmodelled is then the
%   filter's own kernel, whose weight falls as the ratio of the two times (Discussion)."
% WHY IT BELONGS HERE AND NOT ONLY IN THE DISCUSSION. Methods already states the simulator's
% acquisition model twice in this paragraph, and this is the section a reader consults to decide
% whether their data can be run through macroir at all. The Discussion prices the approximation;
% Methods should say what to do.
% ONE PROPERTY WORTH NAMING WHEREVER THIS GOES, because it is specific to this family and reads
% as an advantage rather than as a repair: block-averaging produces EXACTLY the observable of
% Eq. obs-avg at any decimation factor, since the likelihood models the integral over the window
% and not a point sample of a band-limited signal, so decimation raises no aliasing question. A
% method built on instantaneous samples cannot say that.
% [OPEN] The bench-side version of the same statement, which is the more useful one for a reader
% designing an experiment rather than salvaging a record: set the filter near the Nyquist
% frequency of the acquisition window you intend to analyse at. Whether Methods carries the
% design advice or only the salvage instruction is Luciano's call.

Ground truth is generated by exact simulation of the continuous-time Markov chain
(CTMC) by uniformization, so the interval average of the conductance is realised
from an actual sampled trajectory rather than from any moment-closure
approximation. Uniformization draws the transition times themselves, so there is no
within-interval discretisation to choose and none to justify. The production scripts
do pass a \texttt{number\_of\_substeps} argument, and on this branch it has no
effect: the simulation parameters are built from the algorithm name alone, and the
substep count they carry is zero.
% src: projects/eLife_2025/ops/local/figure_3_mle.macroir:69-74 (simulation_algorithm="uniformization"; number_of_substeps passed but inert)
% src: src/core/simulate.cpp:41-52, 325-345 (simulation_parameters_from_algorithm takes the name only)
% src: legacy/qmodel_types.h:1624-1627 (uniformization -> Simulation_n_sub_dt(0))
% src: legacy/qmodel.h:8079-8095 (uniformization branch calls uniformization_sample(mt,...,fs); no substep count)
This is what licenses treating the simulation as the reference against which the
Gaussian likelihood approximations are judged: the approximations are the object
under test, and the simulator supplies the exact answer they are compared to. Every
maximum-likelihood grid cell reported here, the least-squares arm included, used
$10^4$ recordings, with one exception: the cell that carries the coverage result
($N_{\mathrm{ch}}=10$, $\widetilde{S} = 0.005$, boundary-conditioned member) was run at
$10^5$, because coverage of a nominal $95\%$ region in six parameters needs a thousand
independent fits and those fits are taken at group size $100$, so $10^5$ recordings are
required to produce them. The time-resolved run used $1{,}000$ recordings; the
% AMENDED 2026-08-04. The blanket "every cell used 10^4" was false and it was the stop point of a
% replication attempt: the Discussion's coverage sentence says "a thousand fits ... at group size
% 100", which needs 10^5 recordings, and Methods elsewhere states that 10^4 gives "1,000 groups of
% 10 and 100 groups of 100". The data settle it: the cell exists on disk as
% figures/data/0ffbda7/figure_3_G_nch_10_nsim_100000_macro_IR_noise_0.05_*.csv, and it is the only
% nsim_100000 cell in the three freeze directories (five files; every other cell is nsim_10000).
% So the ensemble is real and the sentence describing the study was the thing that was wrong.
single-recording illustration (Figure~1) was drawn by the substep sampler instead,
one $20$-channel trace at $250$ substeps per interval, and does not enter the
quantitative comparison.
% src: projects/eLife_2025/ops/local/figure_1.macroir:30 (number_of_substeps=250 and no simulation_algorithm -> substep sampler)
% src: papers/_program/provenance.md#5 (n_sim=10000 per cell); #6 (1000 recordings); #4 (Figure 1: 20 ch, 250 substeps)
% src: projects/eLife_2025/figures/paper_both/figure_4_common.R:41 (filename pattern hard-codes nsim_10000)
% src: projects/eLife_2025/figures/data/0ffbda7 (93 nonlinearsqr files, all nsim_10000, 30 cells)

The simulation seed was set to $0$, which in this engine is the sentinel for a
random seed drawn from \texttt{std::random\_device}, and the resolved seed was
never logged. The deposited ensembles are therefore statistically equivalent to,
but not bit-for-bit reproducible from, a rerun of the same script. Because each
ensemble is large ($1{,}000$ to $10{,}000$ recordings), the reported statistics are
stable; the Monte-Carlo standard here is statistical equivalence, not bit identity.
% src: papers/_program/provenance.md#9 (seed=0 -> random_device; legacy/mcmc.h:37-45)
% src: papers/1_method/decisions/D-0_freeze_and_rerun_scope.md (seed never logged; not bit-reproducible)

\subsection*{The likelihood family and its members}

The members compared here answer one question first, then a second. The first
question is whether the gating fluctuations are modelled at all. Classical nonlinear
least squares answers no: it fits the deterministic mean current and treats every
departure from it as independent measurement error. The macroscopic likelihoods
answer yes, and among them the second question is how much of the acquisition
interval the interval-averaged conductance is conditioned on. Least squares is the
bottom rung of that cost ladder and the comparison anchor of this paper.

Least squares is selected by the family flag \emph{family approximation} $= 2$ and runs
non-recursive; the variance flag is accepted and ignored on this branch. The averaging flag
still applies to it, and it is what separates the two least-squares arms: at $0$ the fit
predicts the deterministic mean current at one instant, the middle of the acquisition window
(\texttt{LSE}, data key \texttt{nonlinearsqr\_g}), and at $1$ it predicts that same mean
averaged over the window (\texttt{ILSE}, data key \texttt{nonlinearsqr}).
% CORRECTED 2026-08-06. The clause read "at the sampling instant", which places the evaluation at
% the end of the window; every av=0 member evaluates at its middle. The occupancy the prediction is
% taken from is p_start advanced by P_half = expm(Q*Delta/2), and the window is closed by a second
% half step, so av=0 is a midpoint rule and not an endpoint sample. Read at legacy/qmodel.h:4489-4505
% and 4531-4538 (t_P = P_half at av=0, P otherwise), with the same statement in the comment at
% qmodel.h:7513-7518 for the least-squares branch. This is also why the two av=0 arms agree with
% their macroscopic partners to 1e-16 rather than being displaced by half an interval. Each arm therefore
shares its mean model exactly with the macroscopic member carrying the same averaging flag,
\texttt{LSE} with \texttt{NR} and \texttt{ILSE} with \texttt{INR}, and differs from it only
in the predictive variance: on the recording of Figure~\ref{fig:ladder} the predicted means of
\texttt{LSE} and \texttt{NR} agree to $9\times10^{-16}$ interval by interval, and those of
\texttt{ILSE} and \texttt{INR} to $2\times10^{-16}$, while the two least-squares arms differ
from each other by $5\times10^{-2}$ on the same intervals. Both arms
sit off the lattice of endpoint and recursion flags that organises the macroscopic members;
the \texttt{I} prefix records the window setting and does not place \texttt{ILSE} on that
lattice.
% src: projects/eLife_2025/ops/slurm/dispatch_figure_3_LSE.sh:144,146 (nonlinearsqr: recursive=false, averaging=1, family=2; nonlinearsqr_g: same with averaging=0)
% src: projects/eLife_2025/ops/local/figure_3_time_LSE.macroir:67-72 (family_approximation = 2; variance flag ignored, [VAR-BLOCK])
% src: projects/eLife_2025/ops/local/figure_3_time.macroir:115-155 (the two blocks in one script, one ensemble)
% src: the three mean-model numbers were MEASURED here, not taken from the script's note, by reading
% y_mean out of figures/data/figure_1_likelihood_diagnostic_{LSE_av0,NR,LSE,INR}.csv and matching on
% (sample_index, step_start), 12 intervals: max|LSE-NR| = 8.882e-16, max|ILSE-INR| = 2.220e-16,
% max|LSE-ILSE| = 5.469e-02. Checked 2026-08-06.

Every cell of the design grid is run with \texttt{ILSE}: the un-averaged arm was never
dispatched on the grid, and there is no \texttt{nonlinearsqr\_g} output under
\texttt{figures/data}. \texttt{LSE} is measured at the single cell of the time-resolved
figure, from the same script that produces the ensemble the other members score, and drawn
on the illustrative recording of Figure~\ref{fig:ladder}.
% src: projects/eLife_2025/figures/data (0ffbda7: 99 nonlinearsqr files, 82b956f: 8; zero nonlinearsqr_g in any directory, checked 2026-08-06)
% PROVENANCE, load-bearing for anyone reproducing the av=0 arm: before 2026-08-05 the av=0 branch of
% nonlinearsqr propagated the threaded occupancy by P_half and evaluated the mean before advancing,
% so the member fell half a step further behind on every interval. Fixed at legacy/qmodel.h:7490.
% Both LSE dumps used here (figure_1_likelihood_diagnostic_LSE_av0.csv, figure_3_time_dlik_LSE_av0.csv)
% postdate the fix; any earlier dump of that arm is wrong. Source: ops/local/figure_3_time.macroir:141-147.
The macroscopic members all take \emph{family approximation} $= 0$ and are one filter
run under flags that set how the interval-averaged conductance is conditioned. The
\emph{recursive} flag sets whether the occupancy covariance is propagated between
intervals. The \emph{averaging} flag $\mathrm{av}\in\{0,1,2\}$ counts the number of
interval endpoints the hidden state is conditioned on, and it is the program's name for the
window axis of the Theory section: $0$ an instant, $1$ the state at the window's start, $2$
the states at both of its ends. A third flag, the
\emph{variance form}, distinguishes the two members that condition on the interval
start: MR predicts the interval variance with the total per-start-state conductance
variance, while VR removes the between-endpoint part, keeping MR's mean and gain.
% src: projects/eLife_2025/ops/slurm/dispatch_figure_3_G.sh:196-201 (algorithm -> recursive, averaging, variance_form flags)

\begin{table}[htbp]
\begin{fullwidth}
\centering
\begin{tabularx}{\linewidth}{@{}llccl>{\raggedright\arraybackslash}X@{}}
\toprule
Member & Data key & Recursive & $\mathrm{av}$ & Variance & Conditioning of the conductance \\
\midrule
LSE & \texttt{nonlinearsqr\_g} & no & 0 & n/a      & none; the deterministic mean current at the sampling instant \\
ILSE & \texttt{nonlinearsqr} & no  & 1 & n/a      & none; the deterministic mean current averaged over the interval \\
NR  & \texttt{macro\_NR}    & no  & 0 & n/a      & instantaneous; the acquisition averaging is ignored \\
INR & \texttt{macro\_INR}   & no  & 1 & total    & interval mean, without recursion; the endpoints are not separable \\
R   & \texttt{macro\_R}     & yes & 0 & n/a      & instantaneous; the acquisition averaging is ignored \\
MR  & \texttt{macro\_MR}    & yes & 1 & total    & interval mean given the start state \\
VR  & \texttt{macro\_VR}    & yes & 1 & residual & MR's mean and gain, with the between-endpoint variance removed \\
IR  & \texttt{macro\_IR}    & yes & 2 & residual & interval mean given both boundary states; the endpoints enter the gain \\
\bottomrule
\end{tabularx}
\end{fullwidth}
\caption{The members as flags. Both least-squares arms take \emph{family approximation} $= 2$ and every
macroscopic member takes \emph{family approximation} $= 0$. Reading the endpoint flag
as a ladder gives no endpoints (NR, R), one endpoint (MR, VR), then both boundary
states (IR); INR sits outside that ordering, because without a recursive update the
boundary state is unobservable and conditioning on the start state and conditioning on
both ends describe the same method. The variance column reports the form in force rather
than the dispatched flag: IR is dispatched with the total flag, but with both endpoints
conditioned on the residual form is used regardless of it, because the between-endpoint
spread is already carried by the boundary term. % src: legacy/qmodel.h:4557 (averaging==2 || variance_form==residual)
Even so, taken from the same state MR and IR assign the interval average the same total
variance, by two decompositions that rearrange into each other (Appendix~\ref{app:members}), so what
separates them is the gain; because that variance divides the gain, the states they carry
part company at the first gated update and the variances they report along a recording part
with them. VR keeps MR's mean and boundary-free gain and replaces the total per-start-state
conductance variance with the residual one, which is not the same thing as carrying IR's
variance.
% CORRECTED 2026-08-06. The sentence read "MR and IR predict the same total interval variance and
% differ only in the gain", with no "from the same state". As written it is a claim about a
% recording, where it is false: on the figure-1 dumps the two agree to the last printed digit while
% they still share a prior and separate at the first update after that. The identity is between the
% two maps from a prior to a predicted variance. Established in
% papers/1_method/figures_build_plan.md:199-220 (with the numbers), stated correctly in
% 05_discussion.tex:101 ("from the same state"), derived in Appendix 1, and re-verified against
% legacy/qmodel.h:4568-4619 on 2026-08-06. The VR clause was also rewritten: calling VR "MR's gain
% with the between-endpoint variance removed" invited the reading that VR carries IR's variance,
% which figures_build_plan.md:220 explicitly retires, since IR's TOTAL predicted variance is MR's. The least-squares arms and the macroscopic members are
written by different producers, so their output files carry different prefixes,
\texttt{figure\_3\_LSE\_} and \texttt{figure\_3\_G\_}. INR is the member published as
MacroINR in \citet{moffatt2025bayesian}, where it served as the control against which
the boundary-conditioned member was compared; it is reported here only in the
supplement.}
% src: papers/_program/decisions.md sec.2 (INR = MacroINR, the published control; supplement member 2026-07-29)
% NOTE 2026-07-31: the caption used to end "because over the range measured it is not separable from NR",
% sourced to D-4 sec.5 (envelope 1.23-3.6e4 vs NR 1.21-3.5e4; 0/336 within 15% for both). It stays
% REMOVED, and now for a second reason. Those numbers were measured on macro_NMR, the build missing the
% N*ms interval-variance term; macro_INR ran on 2026-07-31 for the figure_3 time cell, and there the two
% are plainly separable: +14.78 nats of log-likelihood over NR, mean squared standardized residual
% 0.9995 against NR's 1.1526, and an accumulated information ratio of 20.8 against NR's 13.6.
% src: projects/eLife_2025/figures/paper_both/figure_3_supplement_1.html (INR-against-neighbours block)
% What is NOT settled: that re-run covers ONE design cell. The D-4 sentence was a statement about the
% envelope over the whole plane, and restating it in either direction needs the plane re-run on
% macro_INR, which has not been done. Do not carry the one-cell numbers across to it.
\label{tab:family}
\end{table}

Table~\ref{tab:family} lists the members. Reading the endpoint flag as a ladder,
the macroscopic members condition the conductance on progressively more of the
interval: no endpoints, one endpoint, then both boundary states, with the exact
trajectory (what the simulator supplies) as the unreachable top rung.
% src: papers/_program/nomenclature.md (endpoint ladder)
\texttt{IR} is the top implemented rung: it conditions on the pair of channel states at
the two ends of the interval, which we call the boundary state, and marginalises the
interior analytically. It is specified by the construction of the Theory section and is
placed here as one member of the completed grid; an interval-recursive likelihood of this
family was first applied to macroscopic recordings, of the P2X$_2$ receptor
\citep{moffatt2025bayesian}. It coincides with a time-integrated Kalman filter, a
correspondence stated and quantified in the Discussion.
% CORRECTED 2026-08-07. The sentence read "It is the likelihood introduced for the P2X2 receptor
% \citep{moffatt2025bayesian}, here placed as one member of the completed grid". That attributed
% the member's SPECIFICATION to the 2025 paper, and the specification is this manuscript's, in
% Theory. The 2025 paper is the family's first application to macroscopic recordings, which is what
% the sentence now says and is what can be checked. The same repair is made at
% 05_discussion.tex:129, which carried the same attribution in the sentence that concedes the
% Kalman device; the two must not drift apart.
% WHY IT MATTERS BEYOND TIDINESS. The abstract and the Introduction now both cite that paper for
% the ranking result, so a referee will open it. Anything in this manuscript that equates its
% member with what ran there is a claim about another paper's implementation that this one cannot
% support, and the member specified in Theory is the only object whose calibration is reported
% here. Keep the citation as application and lineage; never as specification.
The construction of the interval-averaged mean and variance, the boundary-state
covariance, and the conditioning step that consumes them are derived in
Appendix~\ref{app:derivation}; we do not restate them here, beyond the two objects every
member scores with, the regularizer of Eq.~\ref{eq:pseudocount} below, and the
correspondence between the symbols of that derivation and the names the implementation
carries.

\paragraph{The likelihood.} Every member of the family delivers, per window, a predicted mean
$\bar y^{\mathrm{pred}}_t$ and a predicted variance $\sigma^2_{\mathrm{pred},t}$, and the recorded
value is scored under the Gaussian they define,
\begin{equation}
  \log q_t \;=\; -\tfrac{1}{2}\Bigl[
    \log\bigl(2\pi\sigma^2_{\mathrm{pred},t}\bigr)
    + \frac{\delta_t^2}{\sigma^2_{\mathrm{pred},t}}\Bigr],
  \qquad
  \delta_t=\bar y^{\mathrm{obs}}_t-\bar y^{\mathrm{pred}}_t,
  \label{eq:loglik}
\end{equation}
summed over windows. The recursion enters only through what
$(\bar y^{\mathrm{pred}}_t,\sigma^2_{\mathrm{pred},t})$ are conditioned on: a recursive member
carries $(\bm{\mu}^{\mathrm{post}},\bm{\Sigma}^{\mathrm{post}})$ of
Eqs.~\ref{eq:mu-post} and~\ref{eq:sig-post} into the next window, so Eq.~\ref{eq:loglik} sums to
$\log\prod_t q_t(\bar y_t\mid\bar y_{1:t-1})$; a non-recursive member propagates from the initial
condition by Eqs.~\ref{eq:mu-prop} and~\ref{eq:sig-prop} alone and the same sum is
$\log\prod_t m_t(\bar y_t)$. Nothing else in this section is used by the fit.

\paragraph{The initial condition.} The recursion opens with the occupancy the model supplies and
with the covariance that occupancy implies for one channel,
\begin{equation}
  \bm{\mu}_0 = \bm{\pi},
  \qquad
  \bm{\Sigma}_0 = \mathrm{diag}(\bm{\pi})-\bm{\pi}\bm{\pi}^\top .
  \label{eq:init}
\end{equation}
For the schemes run here $\bm{\pi}$ is the first basis vector: the record opens at zero agonist
with every channel shut, so $\bm{\Sigma}_0=\bm{0}$ and the spread the filter carries is generated
entirely by the propagation of Eq.~\ref{eq:sig-prop}. Eq.~\ref{eq:init} is also the one moment at
which $\bm{\Sigma}$ genuinely is a single channel's covariance; from the first update onward,
conditioning on a shared recorded value correlates the channels and what the recursion carries is
$\mathrm{Cov}(\mathbf{N})/N_{\mathrm{ch}}$, which is no single channel's covariance. Every factor of
$N_{\mathrm{ch}}$ in this section comes from that distinction.
% src: legacy/qmodel.h:6726-6740 (init: P_mean = P_initial, P_Cov = diagpos(P_mean) - XTX(P_mean),
% with the comment stating the one-hot reading), and legacy/models_simple.h:57-58 for the two-state
% scheme, where P_initial is the unit vector on state 0. Read first-hand 2026-08-09. NOT the
% equilibrium: calc_Peq exists (qmodel.h:1333) and is not what init consumes.

\paragraph{Every symbol against the name it carries in the implementation.}

The equations of Appendix~\ref{app:derivation} are written with the objects the implementation forms, and the correspondence
is one to one. It is given because a specification whose symbols cannot be found in the code is not
a specification. $\circ$ is the entrywise product and $\mathbf{1}$ the vector of ones;
$\overline{m}_{i\to j}=\mathbb{E}[\Delta^{-2}A^2_{0\to\Delta}\mid X_0=i,X_\Delta=j]$ is the
pair-conditioned second moment, so that
$\bar{v}_{i\to j}=\overline{m}_{i\to j}-\bar{\gamma}^2_{i\to j}$.

\begin{center}
\begin{tabular}{@{}lll@{}}
\toprule
symbol & in the implementation & what it is \\
\midrule
$P_{i\to j}(\Delta)$ & \texttt{P} & transition over the window \\
$\bar{\gamma}_{i\to j}$ & \texttt{gmean\_ij} & conductance averaged over the window, given the pair \\
$G_{i\to j}=\bar{\gamma}_{i\to j}P_{i\to j}$ & \texttt{gtotal\_ij} & the same, weighted, so that it can be summed \\
$\overline{m}_{i\to j}$ & \texttt{gsqr\_ij} & its second moment, given the pair \\
$(\bar\gamma_0)_{i}$ & \texttt{gmean\_i} & $=\texttt{gtotal\_ij}\cdot\mathbf{1}$, Eq.~\ref{eq:gammabar} \\
$\overline{m}_{i\to\cdot}$ & \texttt{gsqr\_i} & $=\texttt{gtotal\_sqr\_ij}\cdot\mathbf{1}$ \\
$H_{i\to j}$, collapsed & \texttt{(gtotal\_ij $\circ$ gmean\_ij)}$\cdot\mathbf{1}$ & the square collapsed, not the collapse squared \\
$\bar{v}_{i\to j}$ & \texttt{gvar\_ij} & present in the window object, unused by the recursion \\
$\bm{\mu}_0$, $\bm{\Sigma}_0$ & \texttt{P\_mean}, \texttt{P\_Cov} & the state where the window opens \\
$\bm{\Sigma}_0-\mathrm{diag}\,\bm{\mu}_0$ & \texttt{SmD} & the prior with its single-channel part removed \\
$\bm{\gamma}^\top\bm{\Sigma}\bm{\gamma}$ and its tilde & \texttt{gSg} & what the state contributes to the variance \\
$\bm{\mu}^\top\mathbf{w}$ & \texttt{ms} & the state-dependent noise, Eq.~\ref{eq:mr-vr} \\
$\mathbf{g}_{0\to\Delta}$ & \texttt{gS} & state against current \\
$\sigma^2_{0\to\Delta}$ & \texttt{y\_var} & the predictive variance \\
\bottomrule
\end{tabular}
\end{center}
% src, all read first-hand in legacy/qmodel.h: :1746-1774 (the window objects), :4540-4557 (the
% state), :4574-4611 (gSg and ms), :4641-4643 (gS), :4585-4615 (y_var). The two slots of
% Eq. mr-vr are the two lines of the code: gsqr_i - gmean_i o gmean_i for MR, and
% gsqr_i - (gtotal_ij o gmean_ij).1 for IR (qmodel.h:4605-4611), so the pair-resolved variance
% gvar_ij is neither consumed nor differentiated by the recursion.
% CORRECTED 2026-08-06. Both sentences pointed at the Theory section for the Kalman correspondence,
% which the Theory section never made: it stated the update, it verified nothing about Kalman
% filters. The verification (coincidence to about 1e-8 with an integrated-measurement augmented
% Kalman filter) is in 05_discussion.tex:127, sourced to docs/bibliography/MacroIR_prior_art_map.md.
% The pointer was already wrong before the Theory section was de-Kalmaned on the same day; that
% edit only made it unambiguous.

Two reductions are worth stating because they anchor the family to prior work.
Setting the boundary-conditioned conductance variance to zero in the recursive,
no-endpoint member (R) recovers the macroscopic-fluctuation update of
\citet{moffatt2007estimation}, so R is the recursive ancestor already in print.
% src: papers/1_method/archive/section-plans-20260729/methods.md#M6 (R reduces to Moffatt 2007)
The single structural difference between the M and I families is that the
one-endpoint members drop the boundary cross-covariance term
$N_{\mathrm{ch}}\,\widetilde{\bm{\gamma}^\top\bm{\Sigma}\bm{\gamma}}$ that IR retains, the
tilde being the boundary-pair contraction defined at Eq.~\ref{eq:tilde}. Every member is
written out on one template in Appendix~\ref{app:members}.
% src: papers/_program/nomenclature.md (MR drops N*gamma^T Sigma gamma that IR keeps)

Throughout the reported runs the remaining build flags were held fixed:
\emph{taylor variance correction} off, \emph{micro approximation} off, and the
\emph{variance approximation} on.
% src: projects/eLife_2025/ops/local/figure_3_mle.macroir:81-89 (taylor=false, micro=false, variance=true)
The V in VR names the \emph{variance form} flag, set to $1$, and has no connection to
the \emph{taylor variance correction} flag, which is off in every run reported here;
the Taylor-corrected variants MRT and IRT are not part of this paper. Because the
Taylor variance correction is off, the recursive one-endpoint member that runs is
MacroMR and not MacroMRT.
% src: projects/eLife_2025/ops/slurm/dispatch_figure_3_G.sh:200,202 (macro_VR: variance_form=1, taylor=false; macro_IRT: taylor=true, not run here)
% src: papers/1_method/archive/section-plans-20260729/methods.md#M6 (taylor=false -> MacroMR, not MRT)

The filter carries three numerical safeguards, present because the code runs them and
because anyone reimplementing it meets the same constraints on day one. None is
housekeeping: the first fixes what the boundary-conditioned moments are where a transition
cannot occur, and without the other two the recursive members return no value at all at the
smaller channel counts.

The first sits upstream of the filter cycle, in the construction of the boundary-conditioned
moments of Eq.~\ref{eq:gammabar-vbar}. Both are quotients by the transition probability
$P_{i\to j}(\Delta)$, and that denominator is zero wherever the transition cannot happen
within one window, which on this design is not an edge case: at zero agonist
$k_{\mathrm{on}}[A]$ vanishes identically and $P_{C\to O}$ is zero over the whole pre-pulse
segment. The numerator vanishes with it, so the quotient is $0/0$ in exact arithmetic and, in
floating point, a ratio of two quantities no larger than the rounding error of the
reconstruction that produced them. The implementation does not test for that case. It
reads each quotient as a posterior mean and adds a conjugate pseudo-count $\varepsilon$ to
numerator and denominator alike,
\begin{equation}
  \bar{\gamma}_{i\to j}
    = \frac{[\mathbf{F}_1(\Delta)]_{i\to j}/\Delta
            \;+\; \varepsilon\,\tfrac12(\gamma_{i}+\gamma_{j})}
           {P_{i\to j}(\Delta)\;+\;\varepsilon},
  \qquad
  \overline{m}_{i\to j}
    = \frac{2[\mathbf{F}_2(\Delta)]_{i\to j}/\Delta^{2}
            \;+\; \varepsilon\,\tfrac12(\gamma^2_{i}+\gamma^2_{j})}
           {P_{i\to j}(\Delta)\;+\;\varepsilon},
  \label{eq:pseudocount}
\end{equation}
where $\overline{m}_{i\to j}$ is the second moment of the interval-averaged
conductance over the same boundary pair, and the variance follows by subtraction,
$\bar{v}_{i\to j}=\overline{m}_{i\to j}-\bar{\gamma}^{\,2}_{i\to j}$.
The two prior means are the endpoint averages of Eq.~\ref{eq:endpoint-limit}, so the prior the
variance inherits is $\bigl((\gamma_{i}-\gamma_{j})/2\bigr)^2$ and is never imposed
separately. The weight is $\varepsilon=\varepsilon_{\mathrm{mach}}\,\kappa_F(\mathbf{V})$,
machine epsilon times the Frobenius condition number of the matrix of right eigenvectors of
$\mathbf{Q}$, which bounds the floating-point noise the spectral reconstruction can leave in
$\mathbf{P}(\Delta)$. The pseudo-count is therefore inert wherever the transition probability
stands above that noise, takes over where the quotient carries no information, and introduces
no fitted constant. Both choices in it are deliberate. Keying $\varepsilon$ to
$\kappa_F(\mathbf{V})$ ties the regularizer to the conditioning of the decomposition that
produced the noise rather than to any model constant, and taking the prior from the
conductance vector rather than from the marginals of $\mathbf{F}_1$ keeps it clean when
$\mathbf{V}$ is ill conditioned, which is when those marginals carry the same reconstruction
error the pseudo-count exists to absorb and when the prior is most needed.
% MOVED HERE 2026-08-06 (Luciano) from 02_theory.tex, which now carries the endpoint limit and a
% one-sentence pointer. Grouped with the trust coefficient and the binomial floor because all three
% answer the same reimplementer question and were previously scattered.
% src: legacy/qmodel.h:1691-1694 (min_P_prior = eps_mach * get<kappa_V>), 1160-1179
% (gmean_ij_prior = (g_i+g_j)/2 and gsqr_ij_prior = (g_i^2+g_j^2)/2), 1189-1202 (calc_g_ij_bayes:
% posterior = (data + prior*eps)/(P + eps)), 1743-1750 (both applied, then gvar_ij = gsqr_ij -
% gmean_ij^2 by subtraction, which is why the variance prior is not imposed separately), 1066-1092
% (kappa_V = ||V||_F * ||W||_F, computed in calc_eigen). Read 2026-08-06.
% The "prior from the conductances, not from the marginals" argument is the code's own, at
% qmodel.h:1154-1159, and the rejected alternative it replaced (gtotal_ij / sqrt(P^2+eps^2), a
% smooth 1/P that shrinks towards zero conductance without saying so) is in
% theory/macroir/notes/Gmean_ij_gvarij/bayesian_prior_regularization_of_Qdt.md. That alternative is
% out of the manuscript on purpose: it is history, not method.
% NOT claimed, deliberately: what P_{C->O} evaluates to numerically here. It is zero in exact
% arithmetic; whether the reconstruction returns a hard zero or a rounding-scale entry has not been
% printed, so the text says 0/0 in exact arithmetic and bounds the noise, and asserts nothing about
% that particular entry.

The second holds the occupancy mean on the probability simplex. Nothing in the Gaussian
conditioning of Eq.~\ref{eq:mu-post} knows that $\bm{\mu}$ is a probability vector, and an
undamped step can drive its entries negative or above one, so the mean step is scaled
by a trust coefficient $\alpha_\mu\in(0,1]$, the largest scaling that keeps the posterior
on the simplex. The covariance is treated separately and takes its full step, since the
rank-one down-date of Eq.~\ref{eq:sig-post} is positive semidefinite by construction; the
two trusts are not tied to a common factor. The form of the damping matters here in a way
it would not in an ordinary filter. This paper differentiates the log-likelihood with
respect to $\bm{\theta}$, so a coefficient that is merely continuous is not enough: a hard
$\alpha=\min(1,\,c\,\alpha_p)$ leaves a step in $\partial\alpha/\partial\bm{\theta}$ at the
boundary and pumps variance into the score whenever the system sits near it, and a
discontinuous form leaves a delta. The coefficient is therefore taken in the smooth form
$\alpha = \tfrac12\bigl(1 + c\,\alpha_p - \sqrt{(1-c\,\alpha_p)^2+\epsilon^2}\bigr)$, a softened
minimum of $1$ and $c\,\alpha_p$ that is infinitely differentiable in $\bm{\theta}$ everywhere.
Here $\alpha_p$ is the largest scaling that keeps the posterior mean on the simplex,
$c = 0.9$ is a safety factor applied when the constraint binds, leaving a ten per cent margin
against the boundary, and $\epsilon = 10^{-4}$ sets the width of the smoothing band, so the
residual bias of $\alpha$ below the hard minimum is at most $\epsilon/2$, which is three orders
of magnitude smaller than that margin. Reimplementing this with a hard minimum reproduces the
mean current and corrupts every quantity this paper measures.

The third is a pair of constants: a binomial-count floor of $5$, below which the expected
count in a state is too small for the Gaussian treatment of that state and the member falls
back to its non-recursive form, and a minimum occupancy probability of $10^{-12}$.
% EXPANDED 2026-08-04. The previous paragraph named the two safeguards and gave neither the
% rule nor the form, which a replication audit reported as cannot-implement: the four recursive
% members return NaN at N_ch 10 and 100 without them. Source: docs/math/trust_coefficient.md
% (281 lines), which carries the decoupling of the two trusts, the softmin form with its epsilon,
% and the history: a discontinuous alpha gave "catastrophic score variance at high N" and the
% C0-continuous min() form "still pumped variance into the score whenever the system hovered near
% the boundary". That history is why the smoothness is stated as load-bearing rather than as a
% detail. Implementation: legacy/qmodel.h:7273-7310 for the binomial floor demoting a recursive
% member to non-recursive.
% CONSTANTS SUPPLIED 2026-08-04, from docs/math/trust_coefficient.md:66-71 (epsilon = 1e-4 as the
% current default, with the residual bias of at most eps/2 = 5e-5 argued there against the margin)
% and :174-175 (trust_multiplying_factor = 0.9, applied when the trust binds, leaving 10% margin).
% Two independent students stopped here in the replication run of 2026-08-04, both on these two
% numbers, after the paragraph above had told them the safeguard exists and what form it takes.
% src: papers/1_method/archive/section-plans-20260729/methods.md#M11 (simplex + PSD trust-region; Binomial_magical_number=5, min_P=1e-12)
% src: theory/macroir/docs/Macro_IR/macro_ir_variance_inflation_correction.tex

\subsection*{The two least-squares parameter configurations}

Least squares runs in two parameter configurations, and they are different statistical
models. Both select \emph{family approximation} $= 2$, both fit data from the same
simulator, and both minimise the same sum of squares. They differ in which parameters
are free, so no quantity computed under one may be compared against the same quantity
computed under the other. Which parameters are free is independent of the window setting
of the previous section, so the configurations and the arms cross: the six-parameter
configuration is run on both arms and supplies the least-squares pair of
Figure~\ref{fig:time}, and the four-parameter configuration is run on \texttt{ILSE} alone
and supplies every grid cell.

The residual scale is not a free parameter of either and is not given a value. It is
integrated out under a scale-invariant prior, so with $\mathrm{SSE}=\sum_i (y_i-\mu_i(\bm{\theta}))^2$
over the $n$ finite recorded values the objective is the marginal
\begin{equation}
  \log L(\bm{\theta}) \;=\; \log\Gamma\!\bigl(\tfrac{n}{2}\bigr)
    \;-\; \tfrac{n}{2}\log\pi \;-\; \tfrac{n}{2}\log \mathrm{SSE}(\bm{\theta}),
  \label{eq:lse-marginal}
\end{equation}
whose maximiser is the ordinary least-squares one, since only the last term depends on
$\bm{\theta}$. What the marginalisation changes is everything downstream of the point
estimate. The score is the derivative of Eq.~\ref{eq:lse-marginal} and carries the factor
$n/\mathrm{SSE}$ rather than a $1/\sigma^2$ that would have to be supplied, and the
information the fit reports is
\begin{equation}
  \mathbf{H}_{\mathrm{LSE}} \;=\; \frac{n}{\mathrm{SSE}}\,
    \sum_i \frac{\partial\mu_i}{\partial\bm{\theta}}\,
           \frac{\partial\mu_i}{\partial\bm{\theta}}^{\!\top},
  \label{eq:lse-fisher}
\end{equation}
the Gauss-Newton curvature with $\hat\sigma^2=\mathrm{SSE}/n$ entering as a plug-in
prefactor and not as a coordinate. There is therefore no $\partial\sigma^2/\partial\bm{\theta}$
term to keep or drop: the scale direction is not in the parameter vector at all, which is
also why the classical arm carries no score in the noise direction and why its mean squared
standardized residual is pinned at one by construction.
% ADDED 2026-08-04. A replication attempt stopped exactly here: "The first thing I could not do
% was give the constant sigma^2 a number. The objective's argmin does not need it, but the score,
% the reported Fisher, the reported interval, the sandwich and the whole 13-to-15-fold LSE
% distortion all [do]." The paper never said the scale is marginalised, so a reader had no way to
% form any of them. Source: legacy/qmodel.h:7382-7398 (the driver's own header, naming the
% Jeffreys marginal and the (n/SSE) J^T J form) and :7566-7592 (sigma2_inv = n / SSE_v, and the
% note that the prefactor uses primitive(SSE) so the AD chain does not differentiate through it).
% The lineage of the marginalised form is moffatt2007responses (JGP).
% OPEN, and it is a real question this paragraph does not settle: the six-parameter configuration
% is described elsewhere as keeping the noise level FREE, which is incompatible with marginalising
% the scale. Either that configuration runs through a different path, or the description is loose.
% Check against the run configs before submission and make the two paragraphs agree.

The six-parameter configuration feeds the time-resolved figure (Figure~\ref{fig:time}). Both
arms of it are produced by \texttt{ops/local/figure\_3\_time.macroir}, the same script that
simulates the ensemble and scores the macroscopic members, so all eight members of that
figure score one common set of recordings.
% src: projects/eLife_2025/ops/local/figure_3_time.macroir:115-155 (the av=1 and av=0 blocks) ; figures/data/figure_3_time_dlik_*.csv (all eight written 2026-08-05 23:13-23:19 from git ccd26f9-dirty)
All six parameters are free in base-ten logarithmic coordinates:
$k_{\mathrm{on}}$, $k_{\mathrm{off}}$, the unitary current, \texttt{Current\_Noise},
the current baseline and the mean channel number.
% src: projects/eLife_2025/ops/local/figure_3_time.macroir:26-31 (six Log10 entries, none Fixed; one par_tr passed to all eight members). The older single-arm script ops/local/figure_3_time_LSE.macroir:39-47 carries the same six and is superseded as the producer.
The script's own header records why. The time-resolved figures accumulate the
per-interval Fisher information and never invert it, so the flat directions of the
least-squares fit are harmless there; and keeping six free preserves the
parameter-index alignment with the macroscopic dumps that the figure code joins
against, whereas fixing two would emit four indices against the macroscopic files' six
and misalign them without any visible error. Two flat directions are then visible
rather than hidden. \texttt{Current\_Noise} has an identically zero score row and an
identically zero Fisher row, because the least-squares mean carries no noise term; and
the unitary current and the mean channel number form a rank-one degenerate block along
the amplitude ridge $N_{\mathrm{ch}}\,i$, so the unitary current is free in this
configuration without being separately identified.
% src: projects/eLife_2025/ops/local/figure_3_time_LSE.macroir:9-18 (rationale, param_index alignment, zero noise row, N*g ridge)

The four-parameter configuration is \texttt{ops/local/figure\_4\_LSE.macroir}, which
produces the grid cells, together with \texttt{ops/local/figure\_3\_mle\_LSE.macroir},
which is the same model run with the per-group estimate cloud retained for the few
cells that need one. Both fix the unitary current at $1$ and \texttt{Current\_Noise} at
the cell value, leaving four free parameters: $k_{\mathrm{on}}$, $k_{\mathrm{off}}$,
the current baseline and the mean channel number. Fixing those two is forced rather
than chosen: they are the flat directions of the previous paragraph, and with them free
the parameter Fisher information is singular, so no maximum-likelihood covariance
exists. Both fixed values are the values the recordings were simulated from. The script
builds one parameter vector and passes it both to the simulator and to the fit as its
reference point, so the fixed value and the truth are the same number by construction.
% src: projects/eLife_2025/ops/local/figure_3_mle_LSE.macroir:72-77 and :108-109 (one par_tr, used as the simulated truth and as theta_reference)
Least squares is therefore not handicapped in this configuration but helped: two of the
six parameters are supplied to it exactly right, while the likelihood members estimate
all six. Where it is nonetheless the arm whose reported uncertainty is furthest from
the delivered one, that failure is not a lack of information about those two. This is
the configuration behind every \texttt{ILSE} cell of Figures~4 and~6.
% src: projects/eLife_2025/ops/local/figure_4_LSE.macroir:54-55 (unitary_current Fixed, Current_Noise Fixed)
% src: projects/eLife_2025/ops/local/figure_3_mle_LSE.macroir:54-55 (same two Fixed; Stage 1 cloud retained)
Its parameter indices are $0 = k_{\mathrm{on}}$, $1 = k_{\mathrm{off}}$,
$2 = $ current baseline, $3 = $ mean channel number. There is no unitary-current index
in these files, and index~2 is the baseline.
% src: projects/eLife_2025/figures/data/0ffbda7 (figure_3_LSE_* battery files: param_index 0..3)

The consequence is carried into the Results. A score, a Fisher information, a
distortion matrix or a log-likelihood computed under six free parameters is a
different object from the one computed under four, and the two distortion matrices
differ in dimension, $6\times 6$ against $4\times 4$, so any scalar that accumulates
over eigenvalues is not comparable between them. Every statement about least squares
in the Results is scoped to one configuration and names it. The statement that least
squares cannot deliver the unitary current is a statement about the four-parameter
configuration.

\subsection*{Parameter estimation and the two anchors}

Inference in each grid cell is two-stage. First, the ensemble is partitioned into
groups of consecutive recordings and a separate maximum-likelihood estimate (MLE) is
computed per group, producing a cloud of estimates whose scatter is the empirical
parameter covariance. The partition is by position and any remainder is dropped,
which costs nothing here because the recordings are independent and identically
distributed and $10^4$ divides by both group sizes used.
% src: src/core/likelihood.cpp:4101-4113 (contiguous blocks, idx = g*gsize+k; remainder dropped with a warning)
The optimiser is Gauss--Newton with Levenberg damping, warm-started at the
simulation truth. The script sets the damping schedule (initial damping $10$, growth
factor $3$, maximum damping $10^{10}$) and three limits: at most $100$ iterations, a
relative gradient tolerance of $10^{-6}$ and a value tolerance of $10^{-10}$.
% src: projects/eLife_2025/ops/local/figure_3_mle.macroir:93-95 (lambda_kickoff=10, lambda_factor=3, lambda_max=1e10, grad_rtol=1e-6, dvalue_tol=1e-10; max_iter=gn_max_iter is injected)
% src: projects/eLife_2025/ops/slurm/dispatch_figure_3.sh:107 (GN_MAX_ITER default 100 -> injected as gn_max_iter)
What stops most fits, however, is a fourth criterion the script cannot reach. The
optimiser tests the Newton decrement $\tfrac12\,g^{\mathsf T}H^{-1}g$ first,
requiring it below $10^{-8}$ for the first ten iterations and below $10^{-4}$ after
that; the relaxation is there because past that point a flat likelihood is being
chased below its own noise floor. The three script-set limits act as fallbacks
behind it, and the gradient-norm one is taken relative to the norm of the warm start
rather than of the current $\theta$, so that a parameter running away along a flat
direction cannot inflate its own tolerance and be recorded as converged. The
Newton-decrement constants are compiled in.
% src: legacy/gauss_newton.h:78-111 (newton_dec2_tol 1e-8, newton_dec2_tol_relax 1e-4, tight_iters 10; not exposed to the DSL)
% src: legacy/gauss_newton.h:237-259 (Newton decrement primary; grad-norm fallback measured against the warm-start norm)
Two group sizes were used, $10$ and $100$; with $10{,}000$ recordings these give
$1{,}000$ groups of $10$ and $100$ groups of $100$.
% src: papers/1_method/archive/section-plans-20260729/methods.md#M7 (group_size 10 and 100; 1000 groups of 10, 100 of 100)
Second, a single joint MLE is computed over all replicates at once (one group of
size $n_{\mathrm{sim}}$), giving the pooled estimate $\theta_{\mathrm{pool}}$; its
near-zero gradient certifies that it is a stationary point of the pooled likelihood.
% src: projects/eLife_2025/ops/local/figure_3_mle.macroir:131-148 (joint MLE = theta_pool)

A group whose refit fails outright is dropped from the cloud, so a cloud can hold
fewer estimates than the partition has groups, and its size is conditional on
convergence. Every fit that did complete carries its convergence status, decided by
the criteria above, into the per-group run log. The figures in the body do not
filter on that status; the supplements that restrict themselves to converged fits
say so.
% src: src/core/likelihood.cpp:4150-4154, 4198-4209 (failed refits dropped, not retained as not-a-number)
% src: src/core/likelihood.cpp:4326-4335 (Convergence_Status_Code per group, written to <path>_runs.csv)
% src: papers/_program/provenance.md:138 (figure_2.Rmd applies no convergence filter to the cloud)
% src: projects/eLife_2025/figures/paper_both/figure_4_supplement_coverage.Rmd:44,105 (supplement filters on Convergence_Status_Code)

Warm-starting the per-group fits at the truth is deliberate. It isolates the error
of the likelihood from the error of the optimiser, which is what this paper studies;
it would be the wrong choice for a study of an estimator, which this is not.
% src: papers/1_method/archive/section-plans-20260729/methods.md#M7 (warm-start caveat: study of the likelihood, not the estimator)

The two anchors, the simulation truth $\theta_{\mathrm{sim}}$ and the pooled fit
$\theta_{\mathrm{pool}}$, are distinct points because the Gaussian macroscopic
likelihood is misspecified with respect to the exact simulator. Their difference
$\theta_{\mathrm{pool}} - \theta_{\mathrm{sim}}$ is the misspecification bias. Every
diagnostic below is computed at both anchors, and the two answer different
questions: at $\theta_{\mathrm{pool}}$ the score is zero by construction so bias
washes out, whereas at $\theta_{\mathrm{sim}}$ the bias is visible but the Fisher
information can be indefinite away from the maximum. Each figure states which anchor
it uses.
% src: papers/1_method/archive/section-plans-20260729/methods.md#M7 (two anchors; state the anchor per figure)
% src: papers/_program/provenance.md#5 (theta_pool score ~ 0; theta_sim Fisher can be indefinite)

\subsection*{Diagnostics and the Gaussian-Fisher anchor}

Faithfulness of each approximate likelihood to the simulator is measured with
likelihood-level diagnostics: standardized residuals, the score (the gradient of
the log-likelihood), and an information-distortion matrix that compares the
covariance $J$ of the score across replicates with the model's own Fisher
information. The definitions, sign conventions, and thresholds of these
diagnostics are owned by the Diagnostics section and its metric registry; here we
describe only how they were computed.
% src: papers/1_method/archive/03_metrics_diagnostics.md (diagnostic definitions owned there)

Every diagnostic reported in this paper is anchored on the analytic Gaussian Fisher
information $\bar G = \sum_t \mathrm{GFI}_t$, summed over intervals. This is the
anchor the Diagnostics section writes $\mathbf{H}$, that is
$\mathbf{H} = \bar G = \sum_t \mathrm{GFI}_t$. Its per-interval term is
\[
  \mathrm{GFI}_t = \frac{(\partial_\theta \bar y_t)(\partial_\theta \bar y_t)^{\mathsf T}}{\sigma_t^2}
  + \frac{(\partial_\theta \sigma_t^2)(\partial_\theta \sigma_t^2)^{\mathsf T}}{2\,\sigma_t^4},
\]
which follows from the analytic derivative alone and needs no extra likelihood
passes. It is positive semidefinite by construction, so no anchor used here can be
indefinite.
% src: papers/_program/provenance.md#7 (Gaussian Fisher G_b = sum_t GFI_t; legacy/qmodel.h:6096-6098)
% src: projects/eLife_2025/ops/local/figure_3_mle_LSE.macroir:152-168 (Fim-less gaussian battery; gidm/gfc/gdcc/gdib)

A second Fisher construction exists in the codebase and was used in an earlier
battery: a central difference of the analytic score, with per-coordinate step
$h_i = h_{\mathrm{rel}}\max(|\theta_i|,1)$ at $h_{\mathrm{rel}} = 10^{-5}$,
symmetrised as $(F_{\mathrm{num}} + F_{\mathrm{num}}^{\mathsf T})/2$. Because the score
itself is obtained by automatic differentiation, this is a single numerical derivative
rather than two. That battery (data directory \texttt{433ed13}) is retained in the
deposit as the demonstration that the two constructions target the same information,
an agreement the time-resolved figures exhibit per interval. No number reported in this
paper is taken from it.
% src: projects/eLife_2025/ops/local/figure_3_mle.macroir:151-157 (numerical Fisher; h_rel injected)
% src: projects/eLife_2025/ops/slurm/dispatch_figure_3.sh:98 (H_RELS default 1e-5)
% src: papers/_program/provenance.md#5 (433ed13 = numerical anchor, superseded for the body)
% src: papers/1_method/decisions/D-0_freeze_and_rerun_scope.md (body moved to the Gaussian anchor)

Comparing the score covariance $J$ against the Fisher information is the sandwich, or
information-matrix-equality, diagnostic \citep{white1982maximum}: equality
$\mathbf{H} = J$ holds only when the likelihood is correctly specified, and the extent
to which the two differ is what the distortion matrix quantifies. The comparison has a
further premise on the least-squares branch, where $\bar G$ presumes homoscedastic
residuals; the Diagnostics section states that premise where the distortion is defined.

Uncertainty on every reported statistic comes from a nonparametric bootstrap over
groups, reported as percentile intervals at the $2.5$, $16$, $50$, $84$ and
$97.5\%$ quantiles. Despite the internal column names, no probit transform is
applied; the intervals are ordinary empirical percentiles, verified equal to the
R quantile of type~$6$.
% src: papers/1_method/archive/section-plans-20260729/methods.md#M8 (percentile bootstrap over groups; NOT probit; quantiles listed)
% src: projects/eLife_2025/ops/local/figure_3_mle.macroir:111 (probit_cis = [0.025,0.16,0.5,0.84,0.975])
The diagnostic battery used $100$ bootstrap replicates and a maximum lag of $10$ for
the score autocorrelation, and the empirical-versus-theoretical distortion capstone
used $200$.
% src: projects/eLife_2025/ops/local/figure_3_mle_LSE.macroir:158-166 (n_boostrap_samples 100, max_lag 10) ; :173-176 (capstone n_bootstrap 200)
The per-group estimate cloud is deposited raw. The program does run a group
bootstrap of it, with a floor of $10$ groups below which it falls back to a single
unresampled pass, but that result is held in memory and never written: the cloud's
only output file is the per-group run log, one row per group and component. Every
interval quoted from the cloud in this paper is therefore computed downstream in the
figure code, from the raw estimates.
% src: include/macrodr/cmd/likelihood.h:1179-1188 (write_csv on the cloud emits <path>_runs.csv and nothing else)
% src: src/core/likelihood.cpp:4278-4295 (group bootstrap; min_groups_for_bootstrap falls back to one unresampled pass)
% src: projects/eLife_2025/ops/slurm/dispatch_figure_3.sh:106 (MIN_GROUPS=10)

\subsection*{Regime grid and reduction}

The acquisition interval was swept over seven levels of $\widetilde{\Delta}$, namely
$1$, $0.5$, $0.2$, $0.1$, $0.05$, $0.02$ and $0.01$, realised as
$500$, $250$, $100$, $50$, $25$, $10$ and $5$ raw samples averaged per interval,
that is step durations of $10$, $5$, $2$, $1$, $0.5$, $0.2$ and $0.1\;\mathrm{ms}$,
giving $10$, $20$, $50$, $100$, $200$, $500$ and $1{,}000$ steps in the record.
% src: projects/eLife_2025/ops/slurm/dispatch_figure_3.sh:184-190 (interval_in_tau labels; n_samp; step counts)
The total record duration is held fixed at $100\;\mathrm{ms}$ across the sweep, so
the interval is varied without varying the amount of data. Within each record the
pre-pulse, agonist and post-pulse segments hold a fixed ratio of $20{:}40{:}40$
percent of the steps.
% src: projects/eLife_2025/ops/slurm/dispatch_figure_3.sh:185-190 (segment step counts 2:4:4 ... 20:40:40; total 100 ms)

The noise axis needs its conversion stated once, because the quantity the figures
plot and the quantity the engine consumes differ by a fixed factor. The engine
parameter is \texttt{Current\_Noise}, a white-noise power spectral density in units of
$g^2\,\mathrm{s}$, and it reaches the likelihood only through the variance of one
recorded point, $e_t = \mathrm{Current\_Noise}/\Delta_t$, with no correlation time of
its own.
% src: legacy/qmodel.h:3740-3741 (Current_Noise enters only as e = S*fs/n_samp)
% [BESSEL-BOUND] OPEN 2026-08-06. Canonical note at 02_theory.tex, under the Bessel paragraph.
% "WITH NO CORRELATION TIME OF ITS OWN" IS THE PREMISE THE FILTER BREAKS, and this is the one
% place in the manuscript where it is written as a property of the engine parameter rather than
% as a modelling choice, so it reads as a fact about the noise instead of a fact about the
% simulator. It is exactly right for what was run: Current_Noise reaches the likelihood only
% through e_t = Current_Noise/Delta_t and nothing else. On a real rig the anti-alias filter gives
% the instrumental noise a correlation time of about 0.13/f_c, which is what makes e_t stop
% growing as 1/Delta_t once the window falls below it. Same dimensionless group as everywhere
% else in this file's Bessel notes: f_c*Delta.
% CONSEQUENCE FOR READING THE INTERVAL AXIS, which 01_introduction.tex carries the other half of:
% the sweep holds the spectral density fixed and shortens the window, so e_t rises as 1/Delta.
% That is realisable on a bench only down to the filter's response time. Nothing measured changes
% and no number moves; what is at stake is whether a reader takes the shortest cells of the sweep
% as configurations they could set up.
% [OPEN] Half a sentence here, or nothing here and the clause in the Introduction's design-axis
% paragraph. One of the two, not both.
Two design quantities are reported in the natural units of the model, and a tilde marks
them: time is measured in the channel time constant $\tau = 1/k_{\mathrm{off}}$ and current
in the unitary current $i$. The acquisition interval is therefore
$\widetilde{\Delta} = \Delta\,k_{\mathrm{off}} = \Delta/\tau$, and the instrumental noise is
\[
  \widetilde{S} \;=\; \frac{S\,k_{\mathrm{off}}}{i^2},
\]
with $S$ the instrumental noise power spectral density, so $\widetilde{S}$ is the instrumental
noise accumulated over one channel time constant measured against the single-channel
signal power. It is the only dimensionless combination of $S$, $k_{\mathrm{off}}$ and $i$, and
it carries no free constant.
At the run values $k_{\mathrm{off}} = 100\;\mathrm{s}^{-1}$ and
$i = 1$ the reference cell used throughout is $\widetilde{S} = 0.01$, and $\widetilde{S} = 1$ is
% SYMBOL REMOVED 2026-08-06. The line read "$i = g = 1$". That $g$ is defined nowhere in the
% manuscript, plays no part in the definition of S_tilde above it, which uses only S, k_off and i,
% and collides with the gain $\mathbf{g}$ of the Theory section. Dropping it costs the sentence
% nothing: it is stating the run values that enter the formula.
the level at which the instrumental noise over one channel time constant equals the
single-channel signal power. The accent carries one other meaning in this paper, the
boundary-pair contractions of the Theory section, Eq.~\ref{eq:tilde} and
Eq.~\ref{eq:vector-tilde}; those are written over bold symbols and the two uses never meet in
one expression.
% ADDED 2026-08-06 (Luciano). Delta_tilde replaces the product Delta*k_off, which was written out
% at ~11 sites in Results and Discussion and on every design-plane axis, and the tilde now has ONE
% stated law instead of being an unexplained mark on a single symbol. Why Delta_tilde and not
% tau_tilde or I_tilde, both of which were considered: the letter under the tilde is the QUANTITY
% and the unit lives in the law, so tau_tilde would be 1 by that same rule; and I is already the
% identity matrix (02_theory.tex:448, 03_diagnostics.tex:34) and I_t the per-interval Fisher
% information (this file, ~line 610). The disjointness sentence is also stated in the appendix
% subsection "The tilde operator", which otherwise claims the accent for the contraction alone.
% INTERNAL, and deliberately not in the printed text (2026-08-06). The sweep was declared
% by a string the dispatchers pair with the engine value through a hand-written case
% statement, label = 1000 * Current_Noise, and that string is 10 * S_tilde: the 1000 is
% 1/Delta at the reference interval 0.1*tau folded into the divisor, not k_off/i^2. The
% stored strings are left exactly as written, because they tie each output file to the cell
% that produced it, and every conversion happens on DISPLAY. The full account, including how
% the discrepancy was found while placing recording configurations on Figure 6, is
% projects/eLife_2025/NOISE_AXIS_UNITS.md, which is where anyone re-auditing the original
% records should start. Nothing in the simulations depends on it: they consumed
% Current_Noise. It is a naming artefact of ours and the reader has no use for it, so the
% paper quotes S_tilde and never mentions the other scale.
% src: projects/eLife_2025/NOISE_AXIS_UNITS.md ; ops/slurm/dispatch_figure_3_G.sh:209-222

The mean channel number $N_{\mathrm{ch}}$ and the noise level are the other two map
axes. The results reported here are read from three data directories, each named by
the git commit hash of the build that wrote it: \texttt{1c2ae6f}, \texttt{87889e6} and
\texttt{0ffbda7}, searched in that order with the first hit winning. Every cell they
contribute to a body figure is at $n_{\mathrm{sim}} = 10^4$, which the figure code
enforces by matching \texttt{nsim\_10000} in the filename, so a cell run at a smaller
ensemble cannot enter a map silently.
% src: projects/eLife_2025/figures/paper_both/figure_4_common.R:31,41 (DATA_DIRS search path; nsim_10000 in the pattern)
% src: projects/eLife_2025/figures/paper_both/figure_6.Rmd:80 (same three directories)
The noise axis is a fixed ladder of twelve levels, and every cell of every algorithm sits on one
of them:
\begin{equation}
  \widetilde{S} \in \{0.005,\; 0.01,\; 0.02,\; 0.05,\; 0.1,\; 1,\; 10,\; 10^2,\; 10^3,\; 10^4,\; 10^5,\; 10^6\},
  \label{eq:noise-ladder}
\end{equation}
denser below one, where the crossover the map measures sits, and by decades above it. The level
$\widetilde{S} = 1$ is where the instrumental noise over one channel time constant equals the
single-channel signal power, so the ladder is dense on the side where the gating fluctuations still
carry information and coarse on the side where they do not.
% ADDED 2026-08-04. Three of five replication attempts stopped here: Methods gave the RANGE
% ("noise labels from 0.05 to 10^6", in the internal label scale) and never the levels, so the
% sweep could not be instantiated. The twelve printed above are those values divided by ten.
% The twelve values are the ones present on disk, from
% papers/1_method/decisions/recompute/d4_distortion_envelopes.py re-run 2026-08-04, which prints
% "noise levels present: 0.05 0.1 0.2 0.5 1 10 100 1000 10000 100000 1000000 10000000".
Across the three directories, \texttt{IR} covers
$N_{\mathrm{ch}} \in \{5, 10, 20, 50, 100, 200, 500, 1000, 2000, 5000, 10000\}$ and
all twelve levels; NR covers
$N_{\mathrm{ch}} \in \{10, 100, 1000, 10000\}$ and R covers those and additionally
$5$, $20$ and $50$ at $\widetilde{S} = 0.01$, with the highest levels run only at the larger
channel counts and reaching $\widetilde{S} = 10^6$ at $N_{\mathrm{ch}} = 10^4$; MR and
VR cover the same channel numbers at $\widetilde{S} \in \{0.01, 0.1, 1\}$; and the least-squares arm
covers $30$ cells on $N_{\mathrm{ch}} \in \{10, 100, 1000, 10000\}$, with $\widetilde{S}$ from
$0.01$ up to $10^3$, $10^4$, $10^5$ and $10^6$ respectively.
% src: projects/eLife_2025/figures/data/1c2ae6f + 87889e6 + 0ffbda7 (per-algorithm nch and noise sets, nsim_10000 files)
The noise range is deliberately not rectangular: the highest levels were run only
at the larger channel numbers, where the crossover the map measures sits. The figure
code auto-detects which cells exist on disk at render time and prints the detected grid
into the render log, so the cell set behind each figure is recoverable from that log.
% src: projects/eLife_2025/figures/paper_both/figure_4_common.R:52-70 (autodetect_cells; detected grid printed on every render)

Each grid cell is reduced through five stages: the per-group MLE cloud; the pooled
joint fit $\theta_{\mathrm{pool}}$; the score at both anchors, whose
across-replicate covariance is $J$ and whose per-interval Gaussian terms sum to the
anchor $\bar G$; the diagnostic battery at each anchor; and the
empirical-versus-theoretical distortion capstone, anchored at
$\theta_{\mathrm{pool}}$. Each cell writes five comma-separated-value families (the
estimate cloud, the pooled fit, the two anchor batteries, and the capstone). The
scripts number these stages one, two, four, five and six; the gap at three is the
numerical-Fisher stage, which the Gaussian anchor removes and which the numbering
keeps visible. On the grid-production lane the cloud stage and the capstone are
dropped, because no map reads them there; the cells that need a cloud are produced
by the separate script that keeps that stage.
% src: projects/eLife_2025/ops/local/figure_3_mle_LSE.macroir:105-178 (stages 1, 2, 4, 5, 6; five CSV families per cell)
% src: projects/eLife_2025/ops/local/figure_3_mle.macroir:150-157 (stage 3 = numerical Fisher, present only on the superseded lane)
% src: projects/eLife_2025/ops/local/figure_4_LSE.macroir:105-119,168-171 (Stage 1 and Stage 6 removed on the grid lane)

\subsection*{Reproducibility and computational environment}

The engine is \texttt{macro\_dr}, built as C++20, and each analysis is a script in a
small typed domain-specific language run in a single process
\citep{moffatt2025macroir}. Provenance is self-documenting at three levels: on every
invocation the binary snapshots the assembled script and its full argument list into
a run ledger; every analysis output file stamps the build's short git commit hash as
its first row; and the data directories are named by that hash, so results from
different code versions can never overwrite one another and multi-commit provenance is
documented per file.
% src: papers/_program/provenance.md#2 (run ledger; hash-stamped row 1; hash-named data dirs)
% src: papers/1_method/decisions/D-0_freeze_and_rerun_scope.md (each CSV self-stamps its hash; multi-commit provenance accepted)

The production batteries ran under SLURM at $32$ CPUs per task and $48\;\mathrm{GB}$
of memory ($96\;\mathrm{GB}$ for the parameter-recovery lane), with the linear-algebra
back end pinned to a single thread so that OpenMP parallelises the per-simulation and
bootstrap loops. The environment variable
\texttt{MACRODR\_\allowbreak AXIS\_\allowbreak SERIAL=1} is
load-bearing: it serialises the internal grid-combination loop, without which memory
grows with the number of concurrent combinations and the jobs exhaust memory.
% src: papers/_program/provenance.md#2,#9 (SLURM 32 CPU, 48/96 GB, BLAS 1 thread, MACRODR_AXIS_SERIAL=1)
% src: projects/eLife_2025/ops/slurm/dispatch_figure_3.sh:196-205 (CPUS 32, MEM 48G, MACRODR_AXIS_SERIAL=1)

The one reproducibility limit worth stating twice is the simulation seed. Because the
seed was the random sentinel and the resolved value was not recorded, the exact
deposited ensembles cannot be regenerated bit for bit; a rerun yields a statistically
equivalent, numerically different ensemble. All resampling that operates on a fixed
dataset (the bootstraps, the group partitioning) is deterministic and reproducible
given the data. Deposition of the data and code, with pinned versions, is stated in
the Data and code availability section.
% src: papers/_program/provenance.md#9 (seed not reproducible; bootstrap on fixed data is deterministic)
% src: papers/1_method/decisions/D-0_freeze_and_rerun_scope.md (statistical-equivalence standard; deposition in backmatter)
